\prefacesection{Abstract}

Scientific software projects rely on collections of related artifacts,
including requirements specifications, design documents, source code, tests,
and build instructions. These artifacts present overlapping knowledge to
different audiences, but maintaining independently editable representations
introduces unnecessary redundancy and creates opportunities for omissions,
inconsistency, and loss of traceability. This thesis investigates how a
knowledge-centric approach can reduce that maintenance burden while preserving
the distinct forms that software artifacts require.

The work first analyzes artifacts from scientific-computing case studies to
identify recurring structures, shared domain knowledge, and repeated
projections of the same information. These observations inform Drasil, a
framework that represents this knowledge (concepts, definitions, symbols, units, equations,
and their relationships) as reusable, typed chunks. Artifact recipes
select and organize those chunks, while deterministic generators render the
resulting structures as documentation, source code, and supporting files.

Reimplementations of six case studies provide observational evidence that this
single-source approach enforces consistency across generated artifacts,
preserves explicit provenance, and supports reproducible generation. Encoding
knowledge explicitly also exposed missing definitions, inconsistent symbols,
implicit conversions, and undocumented assumptions in the original artifacts.
Related systems reused shared knowledge to produce software-family variants
with targeted changes. These benefits require greater up-front modeling effort
and impose an onboarding cost on new contributors. The thesis therefore
establishes the feasibility and practical value of knowledge-centric artifact
generation while identifying improved authoring support, broader domain
coverage, and controlled empirical validation as priorities for future work.
